% ============================================================
%  SECTION 5: CONCLUSION
% ============================================================
\section{Conclusion}
\label{sec:conclusion}

This paper investigated whether the persistent performance gap of Large Language Models on
Amharic---a morphologically complex, low-resource Semitic language written in the
Ge\textquotesingle ez (Fidel) script---could be materially reduced by two principled but
computationally modest interventions: \emph{Ge\textquotesingle ez script awareness} and
\emph{agentic text processing}.  The experimental evidence collected on a purpose-built corpus
of 5{,}000 annotated Amharic instances answers each of the three research questions decisively.

% ------------------------------------------------------------
\subsection*{Answers to the Research Questions}

\paragraph{RQ1 — Agentic text processing.}
Decomposing Amharic language understanding into an explicit three-step orchestration
pipeline---morphological segmentation, semantic disambiguation, and coreference
resolution---yields a 75\,\% relative improvement in translation BLEU (0.12 $\to$ 0.21) and
a 27.6\,\% relative improvement in morphological accuracy (0.58 $\to$ 0.74) compared to a
standard multilingual transformer baseline.  The gains demonstrate that exposing latent
morphological boundaries as structured signals before encoding allows the LLM to generate
more faithful and linguistically precise Amharic output, confirming that \emph{agentic
orchestration is an effective strategy for enhancing both comprehension and generation in
Amharic LLMs}.

\paragraph{RQ2 — Ge\textquotesingle ez script awareness.}
Integrating a dedicated character-level convolutional branch over a $\approx$500-symbol
Ge\textquotesingle ez vocabulary, fused with subword representations, improved all downstream
tasks without any agentic intervention.  Morphological accuracy rose by 19.0\,\% (0.58
$\to$ 0.69), ambiguity/coreference F1 by 16.4\,\% (0.55 $\to$ 0.64), and translation BLEU
by 33.3\,\% (0.12 $\to$ 0.16) relative to the baseline.  These results confirm that
\emph{generic multilingual tokenizers impose a structural ceiling on Amharic conversational
intelligence} and that script-aware character encoding is a necessary, not optional,
architectural component for Ge\textquotesingle ez-script language modelling.

\paragraph{RQ3 — Convergence of script awareness and agentic processing.}
The full Agentic model, which layers multi-step reasoning on top of the Script-Aware encoder,
achieved the highest scores across all tasks.  Critically, the gains of the Agentic model over
the Script-Aware model are non-redundant: a further +31.3\,\% BLEU improvement, +7.2\,\%
morphological accuracy, and +12.5\,\% ambiguity F1 are attributable solely to agentic
reasoning applied on top of script-faithful encodings.  Ablation analysis confirmed a
monotonic improvement with each successive agentic step.  Together, these results establish
that \emph{Ge\textquotesingle ez script literacy and agentic cognitive reasoning are
complementary and synergistic}: the former provides high-fidelity input representations; the
latter provides structured inference over those representations, and the combination
outperforms either paradigm applied in isolation.

% ------------------------------------------------------------
\subsection*{Key Contributions}

The primary contributions of this work are as follows:

\begin{enumerate}
  \item \textbf{A script-aware neural architecture for Ge\textquotesingle ez.}  The proposed
        character-level CNN fusion module is the first to target the complete Fidel inventory
        ($\approx$500 characters) within a multilingual transformer framework, providing a
        reusable inductive bias for any NLP task on Amharic or related Ge\textquotesingle ez-%
        script languages (Tigrinya, Tigre).

  \item \textbf{A modular agentic processing pipeline.}  The three-step orchestration
        (morphological segmentation $\to$ semantic disambiguation $\to$ coreference resolution)
        imposes no additional gradient-based training cost over the Script-Aware model, making
        it an accessible and reproducible extension for low-resource settings.

  \item \textbf{An annotated Amharic benchmark corpus.}  The 5{,}000-instance dataset,
        spanning four linguistically motivated task types with fine-grained morphological
        annotations, is a self-contained evaluation resource that can support future Amharic
        NLP research.

  \item \textbf{Quantitative evidence of synergy.}  The ablation study provides the first
        controlled evidence that Ge\textquotesingle ez script awareness and agentic
        orchestration amplify each other, motivating a \emph{script-first, then reason}
        design principle for Semitic conversational AI.
\end{enumerate}

% ------------------------------------------------------------
\subsection*{Limitations}

Several limitations bound the current study.  First, the benchmark corpus, while carefully
curated, contains 5{,}000 instances; scaling to larger corpora would strengthen the statistical
robustness of the reported gains.  Second, the agentic pipeline steps (morphological
segmentation, disambiguation, coreference resolution) are implemented as rule-augmented text
transformations in the current prototype; replacing them with dedicated learned sub-models would
likely yield additional performance improvements.  Third, evaluation is confined to three
metrics (BLEU, Accuracy, F1); human evaluation of fluency and naturalness in conversational
interactions would complement these automatic scores.  Finally, the study uses
Multilingual-MiniLM as the backbone; larger-scale pretrained multilingual models may show
different sensitivity to the proposed augmentations.

% ------------------------------------------------------------
\subsection*{Future Work}

Building on these findings, several research directions are identified as high-priority:

\begin{itemize}
  \item \textbf{Differentiable agentic steps.}  Replacing the deterministic pipeline with
        end-to-end differentiable reasoning modules (e.g., chain-of-thought fine-tuning or
        tool-augmented LLMs) could allow the model to learn optimal decomposition strategies
        directly from Amharic data.

  \item \textbf{Reinforcement learning from Amharic human feedback (RLHF-am).}  Collecting
        native-speaker preference annotations for generated responses and applying preference
        optimization (e.g., DPO or PPO) would target naturalness and cultural appropriateness
        dimensions not captured by $n$-gram metrics.

  \item \textbf{Cross-lingual transfer to Ge\textquotesingle ez-family languages.}  The script-%
        aware architecture and corpus construction methodology are directly transferable to
        Tigrinya and Tigre; evaluating zero-shot and few-shot cross-lingual performance is a
        natural next step.

  \item \textbf{Larger backbone integration.}  Applying the proposed augmentations to
        billion-parameter multilingual instruction-tuned LLMs (e.g., mT5-XXL, Aya) would
        test whether the synergistic gains persist at scale and could inform practical
        deployment of Amharic conversational agents.

  \item \textbf{Spoken Amharic extension.}  Integrating automatic speech recognition
        preprocessing that retains Ge\textquotesingle ez script fidelity would extend the
        framework to voice-based conversational interfaces, broadening its real-world impact.
\end{itemize}

% ------------------------------------------------------------
\subsection*{Closing Remarks}

The Amharic-speaking community represents over 40 million people, yet the language remains
severely underserved by state-of-the-art conversational AI systems.  This work demonstrates
that targeted architectural choices---rooted in a deep understanding of Ge\textquotesingle ez
orthography and Amharic morphosyntax---can close a meaningful portion of the performance gap
without requiring massive additional compute or data resources.  The script-first, agentic-%
reasoning framework proposed here provides a principled, reproducible, and extensible
foundation for the next generation of equitable, multilingual conversational intelligence.
